% Options for packages loaded elsewhere
\PassOptionsToPackage{unicode}{hyperref}
\PassOptionsToPackage{hyphens}{url}
\PassOptionsToPackage{table}{xcolor}
%
\documentclass[
  11pt,
]{report}
\usepackage{amsmath,amssymb}
\usepackage{iftex}
\ifPDFTeX
  \usepackage[T1]{fontenc}
  \usepackage[utf8]{inputenc}
  \usepackage{textcomp} % provide euro and other symbols
\else % if luatex or xetex
  % pdfLaTeX-only repository policy: the original unicode-math/fontspec
  % branch is unreachable under pdfLaTeX and is replaced with a guard.
  \PackageError{cpython-suite}{Compile with pdfLaTeX only}{}
\fi
\usepackage{lmodern}
\ifPDFTeX\else
  % xetex/luatex font selection
\fi
% Use upquote if available, for straight quotes in verbatim environments
\IfFileExists{upquote.sty}{\usepackage{upquote}}{}
\IfFileExists{microtype.sty}{% use microtype if available
  \usepackage[]{microtype}
  \UseMicrotypeSet[protrusion]{basicmath} % disable protrusion for tt fonts
}{}
\makeatletter
\@ifundefined{KOMAClassName}{% if non-KOMA class
  \IfFileExists{parskip.sty}{%
    \usepackage{parskip}
  }{% else
    \setlength{\parindent}{0pt}
    \setlength{\parskip}{6pt plus 2pt minus 1pt}}
}{% if KOMA class
  \KOMAoptions{parskip=half}}
\makeatother
\usepackage{xcolor}
\usepackage{listings}
\newcommand{\passthrough}[1]{#1}
\lstset{defaultdialect=[5.3]Lua}
\lstset{defaultdialect=[x86masm]Assembler}
\usepackage{longtable,booktabs,array}
\usepackage{calc} % for calculating minipage widths
% Correct order of tables after \paragraph or \subparagraph
\usepackage{etoolbox}
\makeatletter
\patchcmd\longtable{\par}{\if@noskipsec\mbox{}\fi\par}{}{}
\makeatother
% Allow footnotes in longtable head/foot
\IfFileExists{footnotehyper.sty}{\usepackage{footnotehyper}}{\usepackage{footnote}}
\makesavenoteenv{longtable}
\setlength{\emergencystretch}{3em} % prevent overfull lines
\providecommand{\tightlist}{%
  \setlength{\itemsep}{0pt}\setlength{\parskip}{0pt}}
\setcounter{secnumdepth}{5}
% CPython Engineering Suite shared presentation layer -- pdfLaTeX compatible
\usepackage[T1]{fontenc}
\usepackage{lmodern}
\usepackage{microtype}
\usepackage{geometry}
\geometry{letterpaper,margin=0.91in,headheight=16pt,headsep=19pt}
\usepackage{xcolor} % [table] passed at top of file (option-clash fix)
\definecolor{CPNavy}{HTML}{13283B}
\definecolor{CPBlue}{HTML}{176BA0}
\definecolor{CPGold}{HTML}{D8A73C}
\definecolor{CPMuted}{HTML}{566979}
\definecolor{CPBg}{HTML}{EFF5F8}
\usepackage{titlesec}
\titleformat{\chapter}[display]{\normalfont\sffamily\bfseries\color{CPNavy}}{\large\color{CPBlue}CHAPTER \thechapter}{9pt}{\Huge}[\vspace{6pt}\color{CPGold}\titlerule]
\titleformat{\section}{\Large\bfseries\sffamily\color{CPBlue}}{\thesection}{0.7em}{}
\titleformat{\subsection}{\large\bfseries\sffamily\color{CPNavy}}{\thesubsection}{0.7em}{}
\usepackage{fancyhdr}
\pagestyle{fancy}
\fancyhf{}
\fancyhead[L]{\footnotesize\color{CPMuted}\sffamily CPYTHON ENGINEERING SUITE}
\fancyhead[R]{\footnotesize\color{CPMuted}\sffamily 3.16.0a1+dev SNAPSHOT}
\fancyfoot[C]{\footnotesize\color{CPMuted}\thepage}
\renewcommand{\headrulewidth}{0.3pt}
\AtBeginDocument{\urlstyle{tt}}
\lstset{basicstyle=\ttfamily\small,columns=fullflexible,keepspaces=true,breaklines=true,breakatwhitespace=false,showstringspaces=false,frame=single,framesep=5pt,rulecolor=\color{CPBg},backgroundcolor=\color{CPBg},xleftmargin=3pt,xrightmargin=3pt}
\usepackage{tcolorbox}
\tcbuselibrary{breakable,skins}
\newtcolorbox{sourcebox}[1][]{enhanced,breakable,sharp corners,boxrule=0.4pt,colback=CPBg,colframe=CPBlue,borderline west={3pt}{0pt}{CPBlue},left=10pt,right=8pt,top=6pt,bottom=6pt,#1}
\newtcolorbox{cautionbox}[1][]{enhanced,breakable,sharp corners,boxrule=0.4pt,colback=CPBg,colframe=CPGold,borderline west={3pt}{0pt}{CPGold},left=10pt,right=8pt,top=6pt,bottom=6pt,#1}
\usepackage{tikz}
\usetikzlibrary{arrows.meta,positioning,shapes.geometric,calc}
\usepackage{enumitem}
\setlist{itemsep=2pt,topsep=4pt}
\usepackage{booktabs}
\usepackage{longtable}
\AtBeginEnvironment{longtable}{\small}
\usepackage{array}
\setlength{\parskip}{5pt}
\setlength{\parindent}{0pt}
\setcounter{tocdepth}{2}
\usepackage{hyperref}
\hypersetup{colorlinks=true,linkcolor=CPBlue,urlcolor=CPBlue,citecolor=CPBlue,pdfauthor={CPython Engineering Documentation Suite}}
\makeatletter
\renewcommand{\maketitle}{%
  \begin{titlepage}\pagecolor{CPNavy}\color{white}\sffamily
  \vspace*{0.8in}
  {\color{CPGold}\rule{2.1in}{4pt}}\\[16pt]
  {\Large\bfseries CPYTHON ENGINEERING DOCUMENTATION SUITE}\\[24pt]
  {\Huge\bfseries \@title\par}\vspace{19pt}
  \vfill
  {\normalsize Source baseline: CPython 3.16.0 alpha 1+dev repository snapshot}\par
  {\normalsize Prepared 16 October 2026}\par
  {\small\color{white!70} Source-traced engineering commentary; not an official Python specification.}\vspace{0.65in}
  \end{titlepage}\nopagecolor\color{black}
  }
\makeatother
\ifLuaTeX
  \usepackage{selnolig}  % disable illegal ligatures
\fi
\usepackage{bookmark}
\IfFileExists{xurl.sty}{\usepackage{xurl}}{} % add URL line breaks if available
\urlstyle{same}
\hypersetup{
  pdftitle={Python Language Technical Reference},
  hidelinks,
  pdfcreator={LaTeX via pandoc}}

\title{Python Language Technical Reference}
\usepackage{etoolbox}
\makeatletter
\providecommand{\subtitle}[1]{% add subtitle to \maketitle
  \apptocmd{\@title}{\par {\large #1 \par}}{}{}
}
\makeatother
\subtitle{Volume 01 \textbar{} Syntax, semantics, objects, imports, and
executable contracts}
\author{}
\date{}

\begin{document}
\maketitle

{
\setcounter{tocdepth}{1}
\tableofcontents
}
\chapter{Scope, authority, and reading
strategy}\label{scope-authority-and-reading-strategy}

\begin{sourcebox}
\textbf{Evidence classification.} \textbf{SOURCE} = statements supported by material present in the uploaded snapshot; \textbf{IMPLEMENTATION} = CPython-specific internals, not guaranteed Python semantics; \textbf{ENGINEERING} = original recommendations, examples, checks, and exercises. The source baseline is \textbf{3.16.0 alpha 1+dev}; missing repository file bodies are noted as evidence limitations.
\end{sourcebox}

This reference describes the language-facing behavior of Python using
the repository's \emph{Language Reference} and \emph{Tutorial}. It is an
engineering condensation, not an exhaustive reproduction of the
reference grammar. CPython implementation behavior is explicitly
separated from language-level behavior. A language claim should be
checked against \path{Doc/reference/} before
depending on behavior that might differ among implementations.

\begin{longtable}[]{@{}
  >{\raggedright\arraybackslash}p{(\columnwidth - 4\tabcolsep) * \real{0.3333}}
  >{\raggedright\arraybackslash}p{(\columnwidth - 4\tabcolsep) * \real{0.3333}}
  >{\raggedright\arraybackslash}p{(\columnwidth - 4\tabcolsep) * \real{0.3333}}@{}}
\toprule\noalign{}
\begin{minipage}[b]{\linewidth}\raggedright
Evidence tier
\end{minipage} & \begin{minipage}[b]{\linewidth}\raggedright
Source of authority
\end{minipage} & \begin{minipage}[b]{\linewidth}\raggedright
Consequence for engineering
\end{minipage} \\
\midrule\noalign{}
\endhead
\bottomrule\noalign{}
\endlastfoot
Language-facing documentation &
\path{Doc/reference/executionmodel.rst},
\path{lexical_analysis.rst},
\path{simple_stmts.rst},
\passthrough{\lstinline!import.rst!},
\path{toplevel_components.rst} & Basis for portable
language reasoning \\
Introductory exposition & \path{Doc/tutorial/} &
Explains use, not every edge condition \\
CPython-specific & \path{InternalDocs/},
\path{Include/internal/},
\passthrough{\lstinline!Objects/!} & An implementation may change
without changing Python language behavior \\
Proposed design & Examples, exercises and review rules below &
Recommendations, not normative requirements \\
\end{longtable}

\textbf{Version boundary.} \passthrough{\lstinline!README.rst!} states
\textbf{3.16.0 alpha 1+dev}, also reflected in
\path{Include/patchlevel.h}. Do not interpret this
development snapshot as a final Python 3.16 release specification.
\path{Doc/reference/grammar.rst} says its published
grammar is derived from \path{Grammar/python.gram};
that grammar source is named in the tree but its complete body is absent
from the supplied text export.

\chapter{Lexical model and source
structure}\label{lexical-model-and-source-structure}

\section{Tokens, encoding,
indentation}\label{tokens-encoding-indentation}

A Python source file is processed as a stream of tokens. Physical lines,
logical lines, comments, string literals, explicit and implicit line
joining, and indentation delimit suites. Unlike languages with
brace-delimited blocks, indentation is syntactically significant. Spaces
and tabs cannot be mixed arbitrarily for indentation; choose a
consistent whitespace policy. Source encoding and lexical constraints
belong to \path{Doc/reference/lexical_analysis.rst}.

\begin{lstlisting}[language=Python]
# An indented suite belongs to an if statement.
def classify(value: int) -> str:
    if value < 0:
        return "negative"
    if value == 0:
        return "zero"
    return "positive"

entries = [
    "a",  # implicit line joining inside brackets
    "b",
]
\end{lstlisting}

\textbf{ENGINEERING:} Make formatting deterministic and keep encoding
explicit at external I/O boundaries; avoid inferring behavior from
layout transformed by code-generation tools. Syntax errors are caught
before runtime code execution; treat syntax errors differently from
exceptions raised during execution.

\section{Grammar and parser
expectations}\label{grammar-and-parser-expectations}

The CPython documentation describes a Parsing Expression Grammar (PEG)
parser. Ordered alternatives and the cut operator
\passthrough{\lstinline!\~!} are parser-definition details. They are not
operators accessible to Python programs.
\path{InternalDocs/parser.md} explains PEG choice
semantics; \path{Doc/reference/grammar.rst} notes
that the documented grammar omits parser recovery/code-generation
internals. A syntax addition must be evaluated in terms of language
semantics, AST, tooling, and backward compatibility, not grammar changes
alone.

\chapter{Names, binding, scope, and
execution}\label{names-binding-scope-and-execution}

\section{Blocks and binding}\label{blocks-and-binding}

\path{Doc/reference/executionmodel.rst} defines code
blocks including modules, class definitions, function bodies and
interactive commands. Name binding is determined by the syntactic block,
with global and nonlocal declarations changing binding scope. A variable
referenced before its local binding is available can trigger
\passthrough{\lstinline!UnboundLocalError!}. Closures retain access to
enclosing bindings rather than automatically snapshotting values.

\begin{lstlisting}[language=Python]
def counter(start: int = 0):
    value = start
    def increment() -> int:
        nonlocal value
        value += 1
        return value
    return increment

next_value = counter(10)
assert next_value() == 11
assert next_value() == 12
\end{lstlisting}

A name refers to an object; assignment binds a name, and rebinding is
distinct from mutating an existing object. When an object is mutable,
aliasing matters. Do not use \passthrough{\lstinline!is!} as a
substitute for \passthrough{\lstinline!==!} except for identity
comparisons such as \passthrough{\lstinline!is None!}.

\begin{lstlisting}[language=Python]
first = [1, 2]
second = first             # two bindings to one list
second.append(3)
assert first == [1, 2, 3]
assert first is second
copy_of_first = first.copy()
assert copy_of_first == first and copy_of_first is not first
\end{lstlisting}

\section{Evaluating expressions and truth
values}\label{evaluating-expressions-and-truth-values}

Evaluation order, subscription, slicing, comparisons, boolean operators,
calls, comprehensions and assignment expressions are governed by the
language reference. \passthrough{\lstinline!and!} and
\passthrough{\lstinline!or!} short-circuit and return operands; they are
not guaranteed to return \passthrough{\lstinline!bool!}. Generator
expressions are lazy. A comprehension's local iteration variables do not
leak as ordinary bindings into its surrounding scope.

\begin{lstlisting}[language=Python]
def safe_ratio(numerator: float, denominator: float | None):
    if denominator is None or denominator == 0:
        return None
    return numerator / denominator

chosen = [] or ["fallback"]
assert chosen == ["fallback"]
\end{lstlisting}

\textbf{Review invariant:} Distinguish absence
(\passthrough{\lstinline!None!}), falsity (\passthrough{\lstinline!0!},
\passthrough{\lstinline!False!}, empty sequences), and an error.
Replacing \passthrough{\lstinline!if x is None!} with
\passthrough{\lstinline!if not x!} changes behavior when zero/empty
values are legitimate.

\chapter{Python's object and callable
model}\label{pythons-object-and-callable-model}

\section{Objects, types, protocols}\label{objects-types-protocols}

Values are objects with identity, type and value. User-defined classes
expose data and behavior through attributes; special methods make
protocols such as iteration, context management and comparison work. The
Python data model is broader than any single concrete CPython object
layout: do not conflate \passthrough{\lstinline!PyObject!} fields with
portable Python semantics.

\begin{lstlisting}[language=Python]
from dataclasses import dataclass

@dataclass(frozen=True)
class Point:
    x: float
    y: float

    def distance_squared(self) -> float:
        return self.x**2 + self.y**2

assert Point(3, 4).distance_squared() == 25
\end{lstlisting}

Here \passthrough{\lstinline!dataclasses!} is a standard-library
convenience; it is not syntax. \passthrough{\lstinline!frozen=True!}
creates an API that rejects typical attribute assignment, but does not
recursively freeze arbitrary mutable values stored inside an instance.

\section{Functions, defaults, decorators,
annotations}\label{functions-defaults-decorators-annotations}

Function definitions evaluate default argument expressions when the
function is created, not anew on each invocation. Avoid mutable defaults
when per-call independence is needed. Decorators replace or wrap the
object resulting from definition. Annotations should be treated
according to the snapshot's \path{Doc/reference/} and
\path{Doc/library/annotationlib.rst}, rather than
assuming they cause automatic runtime type enforcement.

\begin{lstlisting}[language=Python]
def collect(item: str, bucket: list[str] | None = None) -> list[str]:
    if bucket is None:
        bucket = []
    bucket.append(item)
    return bucket

assert collect("A") == ["A"]
assert collect("B") == ["B"]
\end{lstlisting}

\section{Iterators, generators, and
coroutines}\label{iterators-generators-and-coroutines}

The iteration protocol drives \passthrough{\lstinline!for!} loops; a
generator suspends and resumes execution rather than constructing its
entire output. \passthrough{\lstinline!yield from!} delegates iteration.
\passthrough{\lstinline!async def!} defines coroutine functions, and
\passthrough{\lstinline!await!} participates in asynchronous control
flow rather than automatically creating an operating-system thread.

\begin{lstlisting}[language=Python]
def read_batches(rows, size: int):
    if size <= 0:
        raise ValueError("size must be positive")
    batch = []
    for row in rows:
        batch.append(row)
        if len(batch) == size:
            yield batch
            batch = []
    if batch:
        yield batch
\end{lstlisting}

\chapter{Statements, exceptions, and resource
control}\label{statements-exceptions-and-resource-control}

\passthrough{\lstinline!try!}, \passthrough{\lstinline!except!},
\passthrough{\lstinline!else!}, \passthrough{\lstinline!finally!},
\passthrough{\lstinline!raise!}, and context management separate normal
control flow, exceptional control flow and cleanup. A context manager's
\passthrough{\lstinline!\_\_exit\_\_!} may suppress an exception
depending on its return value. Exception groups and
\passthrough{\lstinline!except*!} warrant specific treatment in
concurrent workflows. See
\path{Doc/tutorial/errors.rst},
\path{Doc/builtins/exceptions.rst}, and
\path{Doc/reference/simple_stmts.rst}.

\begin{lstlisting}[language=Python]
from pathlib import Path

def load_lines(path: Path) -> list[str]:
    try:
        with path.open("r", encoding="utf-8") as stream:
            return stream.read().splitlines()
    except FileNotFoundError as exc:
        raise ValueError(f"required input missing: {path.name}") from exc
\end{lstlisting}

\textbf{ENGINEERING:} Catch only errors a boundary can actually handle;
do not broadly catch \passthrough{\lstinline!BaseException!} to hide
cancellation or process-exit signals. Use
\passthrough{\lstinline!with!}/\passthrough{\lstinline!async with!} for
predictable lifetime management. Avoid relying on immediate object
finalization; CPython reference-count behavior is not a language-level
resource-cleanup guarantee.

\chapter{Imports, module execution, and packaging
boundaries}\label{imports-module-execution-and-packaging-boundaries}

\section{Import process}\label{import-process}

\path{Doc/reference/import.rst} separates import
statements from the import system, its finders/loaders and module cache.
Importing executes module top-level code when a module first loads under
a given interpreter module-cache context. Import cycles can expose
partly initialized modules. Relative imports depend on package context.

\begin{lstlisting}[language=Python]
# app/__main__.py
from .service import run

if __name__ == "__main__":
    raise SystemExit(run())
\end{lstlisting}

Prefer explicit package boundaries and few side effects at import time.
Avoid naming local files after standard-library modules such as
\passthrough{\lstinline!typing.py!}, \passthrough{\lstinline!json.py!}
or \passthrough{\lstinline!logging.py!} where shadowing would be
confusing. Use \passthrough{\lstinline!python -m package!} to execute a
package entry point in its package context.

\section{Introspection and dynamic
evaluation}\label{introspection-and-dynamic-evaluation}

\passthrough{\lstinline!compile!}, \passthrough{\lstinline!eval!},
\passthrough{\lstinline!exec!}, dynamic attribute access, and
metaclasses offer extensibility but also enlarge the attack surface.
\passthrough{\lstinline!eval!}/\passthrough{\lstinline!exec!} must not
be treated as safe interpreters for untrusted expressions merely by
restricting the globals dictionary. In a deployment review, distinguish
trusted configuration from user-controlled input.

\chapter{Behavioral contracts and
verification}\label{behavioral-contracts-and-verification}

\begin{longtable}[]{@{}
  >{\raggedright\arraybackslash}p{(\columnwidth - 4\tabcolsep) * \real{0.3333}}
  >{\raggedright\arraybackslash}p{(\columnwidth - 4\tabcolsep) * \real{0.3333}}
  >{\raggedright\arraybackslash}p{(\columnwidth - 4\tabcolsep) * \real{0.3333}}@{}}
\toprule\noalign{}
\begin{minipage}[b]{\linewidth}\raggedright
Concern
\end{minipage} & \begin{minipage}[b]{\linewidth}\raggedright
Language contract
\end{minipage} & \begin{minipage}[b]{\linewidth}\raggedright
Practical check
\end{minipage} \\
\midrule\noalign{}
\endhead
\bottomrule\noalign{}
\endlastfoot
Name resolution & Binding and scope rules & Nested function with
\passthrough{\lstinline!nonlocal!} \\
Aliasing & Assignment binds objects & Mutate through two aliases \\
Iteration & Iterators may be one-shot & Reiterate generator and list \\
Cleanup & \passthrough{\lstinline!with!} invokes context manager
protocol & Force an exception within the suite \\
Import & Modules and packages obey import machinery & Resolve relative
vs top-level module \\
Type hints & Annotations are not universally enforced & Call annotated
function with wrong type \\
Async & Coroutine requires scheduling/await & Compare coroutine object
and awaited result \\
\end{longtable}

\section{Exercises and acceptance
criteria}\label{exercises-and-acceptance-criteria}

\begin{enumerate}
\def\labelenumi{\arabic{enumi}.}
\tightlist
\item
  \textbf{Scope laboratory.} Build three nested functions using
  \passthrough{\lstinline!global!}, \passthrough{\lstinline!nonlocal!},
  and local variables. Explain every lookup and bind; demonstrate one
  \passthrough{\lstinline!UnboundLocalError!} deliberately.
\item
  \textbf{Object model laboratory.} Implement a bounded queue with
  \passthrough{\lstinline!\_\_len\_\_!},
  \passthrough{\lstinline!\_\_iter\_\_!}, and context management. Test
  mutability and edge cases.
\item
  \textbf{Import laboratory.} Create
  \path{pkg/__init__.py},
  \path{pkg/__main__.py}, and
  \path{pkg/service.py}; run via
  \passthrough{\lstinline!python -m pkg!}, then inspect a circular
  import failure.
\item
  \textbf{Semantics review.} Compare \passthrough{\lstinline!is!},
  \passthrough{\lstinline!==!}, \passthrough{\lstinline!in!},
  \passthrough{\lstinline!and!}, and \passthrough{\lstinline!or!} on a
  set of values including \passthrough{\lstinline!0!},
  \passthrough{\lstinline!None!}, an empty list, and a custom object.
\item
  \textbf{Exception contract.} Write a function whose cleanup runs when
  parsing fails; demonstrate chaining with
  \path{raise ... from ...}.
\end{enumerate}

\textbf{Completion evidence:} runnable tests, a short table of observed
behavior, and references into \path{Doc/reference/}
or \path{Doc/tutorial/} supporting each claim. Avoid
explaining CPython refcounts as if they were language specification
text.

\chapter{Source concordance and change
control}\label{source-concordance-and-change-control}

\begin{itemize}
\tightlist
\item
  \textbf{SOURCE:} \passthrough{\lstinline!README.rst!};
  \path{Include/patchlevel.h} --- the snapshot
  version.
\item
  \textbf{SOURCE:}
  \path{Doc/reference/executionmodel.rst},
  \path{lexical_analysis.rst},
  \path{simple_stmts.rst},
  \passthrough{\lstinline!import.rst!},
  \path{toplevel_components.rst},
  \passthrough{\lstinline!grammar.rst!}.
\item
  \textbf{SOURCE:} \path{Doc/tutorial/classes.rst},
  \path{controlflow.rst},
  \passthrough{\lstinline!errors.rst!},
  \passthrough{\lstinline!modules.rst!},
  \path{datastructures.rst}.
\item
  \textbf{IMPLEMENTATION:}
  \path{InternalDocs/parser.md},
  \path{InternalDocs/frames.md},
  \path{InternalDocs/interpreter.md} --- explanatory
  context, not portable guarantees.
\item
  \textbf{ENGINEERING:} design rules, code samples, verification matrix
  and laboratories are authored for this suite.
\end{itemize}

Before upgrading a production system from a stable release to the
development snapshot, reproduce tests under both versions; consult
version-specific deprecations and release documentation. A feature that
appears in \passthrough{\lstinline!main!} is not necessarily finalized
for a release.

\end{document}
